\documentclass{article}
\usepackage{amsmath,amssymb,mathtools}
\begin{document}

\section*{Strict follow-up: Higham Ch. 4 Kahan returned-sum Eq. (4.8)}

I am formalizing Higham, \emph{Accuracy and Stability of Numerical
Algorithms}, 2nd ed., Chapter 4, equation (4.8), for Kahan summation.  Please
answer as a mathematician, not as a programmer.  I need a proof-assistant-sized
route, or a precise obstruction.  Do not answer with broad prose only.

\subsection*{Algorithm}

Use the stored correction sign convention
\[
  y=\operatorname{fl}(x+e),\qquad
  s'=\operatorname{fl}(s+y),\qquad
  e'=\operatorname{fl}(\operatorname{fl}(s-s')+y).
\]
The returned value is the stored \(s_n\), not \(s_n+e_n\).  Initial state is
\(s_0=e_0=0\).

For one step, write
\[
 y=(x+e)(1+\delta_y),\quad
 s'= (s+y)(1+\delta_s),\quad
 q=(s-s')(1+\delta_{\rm sub}),\quad
 e'=(q+y)(1+\delta_e),
\]
with \(|\delta_\bullet|\le u\) when the operations are in range.

\subsection*{Coefficient recurrence already formalized}

For a source coefficient pair \((S,E)\), one homogeneous coefficient step has
\[
\begin{aligned}
S' &= A S + B E,\\
E' &= C S + D E,
\end{aligned}
\]
where
\[
A=1+\delta_s,\quad
B=(1+\delta_y)(1+\delta_s),
\]
\[
C=-\delta_s(1+\delta_{\rm sub})(1+\delta_e),\quad
D=-(1+\delta_y)(\delta_{\rm sub}+\delta_s(1+\delta_{\rm sub}))(1+\delta_e).
\]
For a newly injected source coefficient, the initial pair is \((B,D)\).

In total/correction coordinates \(T=S+E,\ E=E\), the homogeneous step is
\[
\begin{aligned}
T' &= \alpha T + \rho E,\\
E' &= C T + \eta E,
\end{aligned}
\]
where
\[
\alpha=A+C,\qquad
\rho=(B+D)-(A+C),\qquad
\eta=D-C.
\]
The returned stored coefficient is \(S=T-E\).

The following pointwise identities/bounds have already been checked:
\[
\eta+\delta_{\rm sub}=O(u^2),\qquad
\rho-\eta-\delta_y=O(u^2),
\]
hence
\[
\rho+\delta_{\rm sub}-\delta_y=O(u^2).
\]
Also, if \(\delta_{\rm sub}=0\) at every correction subtraction, then the
ordinary returned coefficient satisfies the source-shaped bound
\[
|S_i-1|\le 2u+C n u^2.
\]
The compensated-total coefficient \(T_i\) also has such a bound.  The missing
question is whether \(S_i=T_i-E_i\) has it without exact correction subtraction.

\subsection*{Known obstructions already proved locally}

\begin{itemize}
\item A bare local relative-error model is too weak: it need not make the first
step from \((0,0)\) exact, and a small-\(u\) biased abstract model violates any
fixed \(2u+C u^2\) returned-\(s\) cap.
\item Finite representability of the operands alone does not imply exact
correction subtraction.
\item Tail FastTwoSum order, tail Sterbenz, tail Ferguson, and direct tail
finite subtraction are false for arbitrary input order.
\item A triangle estimate \(S=T-E\) is too weak because \(E\) can be first order.
\item An affine residual proof would work if the final retained-correction
residual were \(O(nu^2)\sum_i |x_i|\), but the current input-only residual
majorant is first order even for one exact input.
\end{itemize}

\subsection*{Required answer}

Please provide exactly one of the following, with enough equations to translate
to Lean:

\begin{enumerate}
\item A valid finite-round-to-even, arbitrary-input-order proof that the
ordinary returned \(s_n\) has coefficients
\[
  s_n=\sum_i x_i(1+\mu_i),\qquad |\mu_i|\le 2u+Cnu^2.
\]
State the invariant on \((T,E)\) or \((S,E)\), and show exactly where the
first-order \(\delta_{\rm sub}\) terms cancel in \(S=T-E\), not only in \(T\).

\item A concrete finite-round-to-even counterexample, with a specified tiny
format or symbolic exact real values, showing that ordinary returned \(s_n\)
can require a coefficient first-order term greater than \(2u\) under arbitrary
input order.  It is not enough to give a bare relative-error counterexample;
say explicitly whether the obstruction is only for the abstract model or also
for actual finite round-to-even.

\item A corrected theorem with the weakest natural hidden hypotheses.  For
example, exact correction subtraction, FastTwoSum order/range, Sterbenz,
Ferguson, a guard-digit/no-underflow condition, or a different returned
quantity \(s_n+e_n\).  If this is the answer, state which Higham/Goldberg proof
step uses the hidden hypothesis and give a proof skeleton with lemmas.
\end{enumerate}

If the answer is ``unknown from these assumptions'', say what additional source
fact must be acquired.  Separate proven facts from conjecture.

\end{document}
